\section{Testing}
I test sono stati eseguiti tramite il framework \texttt{JUnit} e sono stati strutturati come segue:
\begin{verbatim}
 src/test/
`-- java
    |-- IntegrationTest
    |   |-- BookIntegrationTest.java
    |   |-- CardIntegrationTest.java
    |   |-- EventIntegrationTest.java
    |   |-- LibraryIntegrationTest.java
    |   |-- LoanIntegrationTest.java
    |   |-- ORM
    |   |   |-- BaseIntegrationTest.java
    |   |   |-- BookDAOIntegrationTest.java
    |   |   |-- CardDAOIntegrationTest.java
    |   |   |-- ConnectionManagerUnitTest.java
    |   |   |-- EventDAOIntegrationTest.java
    |   |   |-- LibraryDAOIntegrationTest.java
    |   |   |-- LoanDAOIntegrationTest.java
    |   |   |-- RoomDAOIntegrationTest.java
    |   |   `-- UserDAOIntegrationTest.java
    |   `-- StudyRoomIntegrationTest.java
    `-- UnitTest
        |-- BusinessLogic
        |   `-- AuthUnitTest.java
        |-- DomainModel
        |   |-- EventRoomUnitTest.java
        |   |-- EventUnitTest.java
        |   |-- LoanUnitTest.java
        |   |-- RoomUnitTest.java
        |   `-- StudyRoomUnitTest.java
        `-- utils
            |-- DBSeeder.java
            `-- DBSeederTest.java
\end{verbatim}
Sono stati divisi nei package \texttt{UnitTest} e \texttt{IntegrationTest}. Nel primo si va a testare la classe in isolamento, facendo il mock del database se necessario, come fatto un \texttt{AuthUnitTest} di cui si riporta un estratto di codice; mentre nel secondo si vanno a chiamare i metodi implementati nei DAO interrogando realmente il database. Per l'elenco di tutti i test case si rimanda alla sezione \hyperref[tabella_test]{Tabella test}.
\begin{code}
\begin{minted}{java}
@BeforeAll
    public static void init() {
        userDAO = mock(UserDAO.class);
        when(userDAO.getUserByEmail("library.user@email.com")).thenReturn(new UserDTO(
                "library.user@email.com",
                "Mario",
                "Rossi",
                role.libraryUser,
                BCrypt.hashpw("password", BCrypt.gensalt())
        ));
        when(userDAO.getUserByEmail("librarian@email.com")).thenReturn(new UserDTO(
                "librarian@email.com",
                "Luca",
                "Bianchi",
                role.librarian,
                BCrypt.hashpw("bibliotecario", BCrypt.gensalt())
        ));
        when(userDAO.getUserByEmail("admin@email.com")).thenReturn(new UserDTO(
                "admin@email.com",
                "Sara",
                "Verdi",
                role.libraryAdministrator,
                BCrypt.hashpw("admin", BCrypt.gensalt())
        ));
        when(userDAO.getUserByEmail("inesistente@email.com")).thenThrow(new UserNotFoundException("Wrong username or password"));
        when(userDAO.updatePassword(eq("library.user@email.com"), anyString())).thenReturn(true);
        when(userDAO.updatePassword(eq("inesistente@email.com"), anyString())).thenReturn(false);
        when(userDAO.insertUser(any())).thenReturn(false);
        when(userDAO.insertUser(refEq(new UserDTO(
                "nuovo@email.com",
                "Gino",
                "Neri",
                role.libraryUser,
                "password"
        )))).thenReturn(true);
    }
    
    @Test
    public void login_invalidPassword_throwsException() {
        AuthService authService = new ConcreteAuthService(userDAO);
        var e = assertThrows(UserNotFoundException.class, () -> authService.login("library.user@email.com", "Password"));
        assertEquals("Wrong username or password", e.getMessage());
    }
\end{minted}
\caption{init e test case in AuthUnitTest}
\end{code}
\noindent Per agevolare l'automatizzazione degli integration test è stata scritta una classe \texttt{DBSeeder} che va a creare uno schema \texttt{test}, istanziato mediante il file \texttt{db.sql} e popolato tramite il file \texttt{data.sql}. Gli script \textit{SQL} sono eseguiti tramite la classe \texttt{ResourceDatabasePopulator} del package \texttt{springframework.jdbc}
\begin{code}
    \begin{minted}{java}
    public class DBSeeder {

    private Connection connection;

    public DBSeeder() throws SQLException{
        this.connection = ConnectionManager.getInstance().getConnection();
    }

    public void createTestSchema() {
        try {
            Statement stmt = connection.createStatement();
            stmt.executeUpdate("CREATE SCHEMA IF NOT EXISTS test AUTHORIZATION CURRENT_USER");
        } catch (SQLException e) {
            e.printStackTrace();
        }
    }
    public void deleteTestSchema(){
        try {
            connection.setAutoCommit(true);
            Statement stmt = connection.createStatement();
            stmt.executeUpdate("DROP SCHEMA IF EXISTS test CASCADE");
            stmt.executeUpdate("SET search_path TO public");
        } catch (SQLException e) {
            e.printStackTrace();
        }
    }

    public void initDatabaseTest(String filenameDDL, String filenameData){
        try {
            Statement stmt = connection.createStatement();
            stmt.execute("SET search_path TO test");
        } catch (SQLException e) {
            e.printStackTrace();
            throw new RuntimeException("Database error. Check if schema 'test' exists");
        }
        ResourceDatabasePopulator populator = new ResourceDatabasePopulator();
        populator.addScripts(
                new FileSystemResource(filenameDDL), //file must be in root directory
                new FileSystemResource(filenameData)
        );
        populator.populate(this.connection);
    }
}
    \end{minted}
\caption{Classe DBSeeder}
\end{code}
\noindent Il setup e il teardown sono comuni a tutti gli integration test; essi sono definiti nella classe \texttt{BaseIntegrationTest} che ogni classe nel package estende. Il setup va a creare lo schema \texttt{test}, crea le tabelle e imposta l'autocommit della connessione a falso; dopo ogni test case viene eseguito il teardown, dove viene fatto il rollback della transazione. In caso di fallimento di questa operazione, viene impostato un flag \texttt{teardownFail} a vero, che il metodo \texttt{teardownChecker} eseguito dopo ogni test case controlla, facendo fallire tutti i test dato lo stato inconsistente del database. Infine il metodo \texttt{connectionCloser} cancella lo schema precedentemente creato e chiude la connessione.
\begin{code}
    \begin{minted}{java}
    private static void setTestSchema() throws SQLException {
        Statement stmt = connDAO.createStatement();
        stmt.execute("SET search_path TO test");
    }

    @BeforeAll
    public static void init() {
        try {
            dbSeeder = new DBSeeder();
        } catch (SQLException e) {
            e.printStackTrace();
            fail("Failed to instantiate DBSeeder. Check connection to db");
        }
        dbSeeder.createTestSchema();
        assertDoesNotThrow(() -> dbSeeder.initDatabaseTest("db.sql", "data.sql"));
        try {
            connDAO = ConnectionManager.getInstance().getConnection();
            setTestSchema();
            connDAO.setAutoCommit(false);
        } catch (SQLException e) {
            e.printStackTrace();
            fail("Fail to get connection");
        }
    }

    @AfterEach
    public void teardown() {
        try {
            connDAO.rollback();
        } catch (SQLException e) {
            teardownFail = true;
            e.printStackTrace();
        }
    }

    @BeforeEach
    public void teardownChecker() {
        assumeFalse(teardownFail, "Aborting test. DB state is inconsistent");
    }


    @AfterAll
    public static void connectionCloser(){
        try {
            dbSeeder.deleteTestSchema();
            connDAO.close();
        } catch (SQLException e) {
            e.printStackTrace();
        }
    }        
    \end{minted}
    \caption{metodi della classe BaseIntegrationTest}
\end{code}
\subsection{Test coverage}
I risultati di copertura sono stati generati con \texttt{JaCoCo} e sono disponibili per la consultazione su \url{https://salvatore-greco.github.io/SWE/}.\\
È stato raggiunto un buon risultato col \textbf{90\%} di linee di codice coperte e il \textbf{79\%} di branch coperti. Il risultato più basso coi branch è causato da alcune condizioni non testate sui ResultSet nei DAO.

\subsection{Unit Test}
In questo package sono presenti gli unit test di alcune classi del package \texttt{DomainModel}, che implementano metodi effettivamente da testare, di \texttt{ConcreteAuthService} del package \texttt{BusinessLogic} e della classe \texttt{DBSeeder}.

